\section{Attention 与 Transformer：信息怎样被选择和组合}

前一章回答“为什么会走到 Transformer”，本章回答“它具体做了什么”。课堂从 seq2seq 进入 Q/K/V，再区分 self-attention、multi-head attention 与 cross-attention。理解这些差异，才能在后面判断 decoder-only LLM、encoder--decoder、retrieval 与 agent memory 的接口位置。

\subsection{Seq2seq 的瓶颈：单个 state 承担太多}

早期 sequence-to-sequence（序列到序列）模型用 encoder 把输入压进隐藏状态，再由 decoder 逐步生成。长度增加时，单个向量必须同时保存词义、顺序与远距离依赖，信息瓶颈会越来越明显。下面两页先展示 recurrent path，再展示 attention 增加的直接读通路；比较重点是信息从哪里流向 decoder。

\lecturefigure{slide-019.jpg}{RNN seq2seq 通过 recurrent state 传递序列信息}{官方 CS25 V4 slide deck，第 19 页}

\begin{importantbox}{Recurrent update}
最简 RNN 可以写成
\[
h_t=\phi(W_xx_t+W_hh_{t-1}+b),\qquad y_t=W_oh_t.
\]
$x_t$ 是第 $t$ 个输入，$h_t$ 是隐藏状态，$W_x,W_h,W_o$ 是可学习矩阵，$\phi$ 是非线性函数。每一步依赖 $h_{t-1}$，因此时间维难以完全并行；长序列还要让梯度穿过很多 recurrent steps。
\end{importantbox}

\lecturefigure{slide-020.jpg}{Attention 允许 decoder 直接访问多个 encoder state}{官方 CS25 V4 slide deck，第 20 页}

\begin{knowledgebox}{读图：Attention 改变的是信息通路}
没有 attention 时，decoder 主要依赖最终 encoder state；加入 attention 后，每个输出位置都能对全部输入位置打分并做加权求和。模型不再被迫把整句压成一个固定瓶颈，而是按当前生成需求重新读取输入。
\end{knowledgebox}

\subsection{Q、K、V：查询、索引与内容}

课堂用 library system 类比 Query、Key、Value。读者提出查询，书目索引承担 key，书中内容承担 value。模型先比较 query 与各 key 的相关性，再用权重组合对应 value。读图时要把“检索到哪本书”和“从书中取出什么”分开，这正对应 key matching 与 value aggregation 两个步骤。

\lecturefigure{slide-021.jpg}{Library analogy 解释 Query、Key 与 Value}{官方 CS25 V4 slide deck，第 21 页}

\begin{importantbox}{Scaled dot-product attention}
令 $Q\in\mathbb{R}^{n_q\times d_k}$、$K\in\mathbb{R}^{n_k\times d_k}$、$V\in\mathbb{R}^{n_k\times d_v}$，attention 为
\[
\operatorname{Attention}(Q,K,V)=\operatorname{softmax}\!\left(\frac{QK^\top}{\sqrt{d_k}}+M\right)V.
\]
$QK^\top$ 给出每个 query 与 key 的匹配分数；$\sqrt{d_k}$ 控制点积尺度；$M$ 是可选 mask；softmax 把分数变成权重；最后对 $V$ 做加权求和。Query 决定“现在要找什么”，Key 决定“我是否相关”，Value 决定“相关时提供什么内容”。
\end{importantbox}

\begin{warningbox}{Attention weight 不是完整解释}
权重展示了某层某头如何混合 value，但输出还经过 value projection、residual、MLP 和后续层。不能仅凭一张 attention map 就断言模型“因为什么做出决定”，更不能把高权重直接解释为因果贡献。
\end{warningbox}

\lecturefigure{slide-022.jpg}{Attention weight 将多个 value 组合成 context vector}{官方 CS25 V4 slide deck，第 22 页}

\begin{lstlisting}[caption={Scaled dot-product attention 的教学版伪代码}]
def attention(query, key, value, mask=None):
    scores = query @ key.transpose(-1, -2)
    scores = scores / sqrt(key.shape[-1])
    if mask is not None:
        scores = scores + mask
    weights = softmax(scores, dim=-1)
    return weights @ value, weights
\end{lstlisting}

\teachervoice{讲者强调这是“soft match”：一个 query 不必只选一个 value，而可以从多个位置取信息。这个性质让 attention 与数据库精确匹配不同，也解释了它为什么能表达模糊、组合式的语言关系。}

\subsection{Self-attention 与 Multi-head attention}

Self-attention（自注意力）让同一序列同时产生 $Q$、$K$、$V$。每个 token 都能根据内容读取其他 token；causal decoder 会用 mask 禁止读取未来位置。Multi-head attention（多头注意力）则让多个子空间并行学习不同的匹配模式。接下来的三页从单头连接、完整 block 到多头可视化逐层展开，先看信息来源，再看 projection 与融合。

\lecturefigure{slide-023.jpg}{Self-attention 在同一序列内部建立内容相关连接}{官方 CS25 V4 slide deck，第 23 页}

\begin{knowledgebox}{背景概念：self-attention 的三种常见 mask}
\begin{tabularx}{\textwidth}{L{0.23\textwidth} X}
\toprule
形式 & 允许的信息流 \\
\midrule
Bidirectional & 每个位置可读取整个输入，典型于 encoder。 \\
Causal & 位置 $t$ 只能读取 $\le t$ 的 token，典型于 autoregressive decoder。 \\
Block/sparse & 只允许局部块、全局 token 或预定义连接，用于降低长序列成本。 \\
\bottomrule
\end{tabularx}
\end{knowledgebox}

\lecturefigure{slide-024.jpg}{Transformer block 将 multi-head attention 与前馈网络堆叠}{官方 CS25 V4 slide deck，第 24 页}

\begin{importantbox}{Multi-head attention}
第 $i$ 个 head 先做独立投影：
\[
\operatorname{head}_i=\operatorname{Attention}(QW_i^Q,KW_i^K,VW_i^V),
\]
再把所有 head 拼接并输出：
\[
\operatorname{MHA}(Q,K,V)=\operatorname{Concat}(\operatorname{head}_1,\ldots,\operatorname{head}_h)W^O.
\]
$h$ 是 head 数，$W_i^Q,W_i^K,W_i^V$ 是每个 head 的投影，$W^O$ 负责融合。多头并不保证每个 head 自动对应清晰语义，但它增加了并行表示子空间。
\end{importantbox}

\lecturefigure{slide-025.jpg}{不同 attention head 可学习不同位置关系}{官方 CS25 V4 slide deck，第 25 页}

\begin{warningbox}{“一个 head 负责一种语法”只是观察，不是接口契约}
某些 head 会呈现位置、指代或句法模式，但训练没有要求 head 永久承担固定角色。剪枝、再训练或换模型后模式可能变化。工程设计不能依赖某个编号 head 必然保存某类知识。
\end{warningbox}

\subsection{Cross-attention：两个序列之间怎样对齐}

Cross-attention（交叉注意力）把 query 与 key/value 放在不同来源。机器翻译 decoder 产生 query，encoder hidden states 提供 key 和 value；检索增强、视觉语言模型和 encoder--decoder 也会使用类似接口。

\lecturefigure{slide-026.jpg}{机器翻译中的 self-attention 与 cross-attention}{官方 CS25 V4 slide deck，第 26 页}

\begin{knowledgebox}{术语表：三种 attention 不要混在一起}
\begin{tabularx}{\textwidth}{L{0.22\textwidth} L{0.25\textwidth} X}
\toprule
类型 & Q/K/V 来源 & 主要用途 \\
\midrule
Self-attention & 同一序列 & 建立序列内部依赖。 \\
Causal self-attention & 同一序列，带未来 mask & 自回归生成。 \\
Cross-attention & Q 来自 decoder/一种模态，K/V 来自 encoder/另一模态 & 条件生成与跨模态、跨序列融合。 \\
\bottomrule
\end{tabularx}
\end{knowledgebox}

\subsection{Transformer 相比 RNN 赢在哪里}

最后一页对比把优势归纳为长距离访问、稳定训练与并行计算。需要补上一个限制：标准 dense self-attention 的 score matrix 是 $n\times n$，序列长度增长会带来二次计算和内存成本。读这张表时应同时比较 dependency path、sequential steps 与 resource cost，避免只记“Transformer 更好”。

\lecturefigure{slide-027.jpg}{RNN 与 Transformer 的训练和依赖建模对比}{官方 CS25 V4 slide deck，第 27 页}

\begin{importantbox}{复杂度与并行性的交换}
对长度 $n$、隐藏维度 $d$ 的序列，dense attention 的主要 score 计算约为 $O(n^2d)$，attention matrix 占 $O(n^2)$；RNN 单步约为 $O(d^2)$，总计算随 $n$ 线性增长，却必须沿时间顺序执行。Transformer 用更多全局配对换取了训练并行性与短的信息路径。
\end{importantbox}

\teachervoice{课堂把 GPU parallelism 视为 Transformer 成功的重要原因。架构不仅要表达得好，还要能匹配硬件吞吐；这也是后续 scale、成本与研究集中化问题的起点。}

\subsection{本章小结}

Attention 通过内容相关加权让模型按需读取信息。自注意力处理序列内部关系，交叉注意力连接不同来源，多头机制提供多个表示子空间。Transformer 由此摆脱 recurrent bottleneck，却引入长序列的二次成本和更复杂的可解释性问题。

